\subsection{Finding New Coordinates from a Linear System}

The geometric picture tells us what must be true. To calculate the coefficients,
write
\[
\mathbf b_1=(a_1,a_2),
\qquad
\mathbf b_2=(b_1,b_2),
\qquad
\mathbf u=(x,y).
\]

\begin{derivationbox}[From a vector equation to scalar equations]
The equation
\[
\mathbf u=c_1\mathbf b_1+c_2\mathbf b_2
\]
means equality of corresponding Cartesian components. Therefore
\[
\boxed{
\begin{aligned}
 a_1c_1+b_1c_2&=x,\\
 a_2c_1+b_2c_2&=y.
\end{aligned}}
\]
This system of two linear equations is the algebraic content of changing from
the Cartesian basis to $\mathcal B$.
\end{derivationbox}

\begin{examplebox}
For
\[
\mathbf b_1=(2,1),
\qquad
\mathbf b_2=(1,2),
\qquad
\mathbf u=(5,4),
\]
we solve
\[
\begin{aligned}
2c_1+c_2&=5,\\
c_1+2c_2&=4.
\end{aligned}
\]
From the first equation,
\[
c_2=5-2c_1.
\]
Substituting into the second gives
\[
c_1+2(5-2c_1)=4,
\]
so
\[
-3c_1=-6,
\qquad
c_1=2.
\]
Then
\[
c_2=5-2\cdot2=1.
\]
Hence
\[
\boxed{[\mathbf u]_{\mathcal B}=(2,1)}.
\]
The calculation agrees with the oblique parallelogram.
\end{examplebox}
